\subsection{单生成超越扩张的情形}

引理 1. 设 K 为域，v 为 K 上的赋值，Γ 为其阶群，Γ' 是一个全序的

群，包含 \(\Gamma\)，且 \(\xi\) 是 \(\Gamma'\) 的一个元素。存在唯一的赋值 \(w\)，定义在 \(\mathbf{K}(\mathbf{X})\) 上，使得对于 \(\mathbf{P} = \sum_{j} a_j \mathbf{X}^j (a_j \in \mathbf{K})\)，有 \(w(\mathbf{P}) = \inf_j(v(a_j) + j\xi)\)。

根据 § 3 第 2 节的命题 4，只需证明公式

\[
w \left(\sum_ {j} a _ {j} \mathbf {X} ^ {j}\right) = \inf _ {j} (v (a _ {j}) + j \xi)\tag{1}
\]

在环 K{[}X{]} 上定义一个赋值。由于

\[
v \left(a _ {j} + b _ {j}\right) + j \xi \geqslant \inf (v \left(a _ {j}\right), v \left(b _ {j}\right)) + j \xi = \inf (v \left(a _ {j}\right) + j \xi , v \left(b _ {j}\right) + j \xi),
\]

由此可得

\[
w (\mathbf {P} + \mathbf {Q}) \geqslant \inf (w (\mathbf {P}), w (\mathbf {Q}))\tag{2}
\]

对于 \(\mathbf{P},\mathbf{Q}\) 属于 \(\mathbf{K}[\mathbf{X}]\)，当 \(w(\mathbf{P})\neq w(\mathbf{Q})\) 时取等号。我们证明

\[
w (\mathrm{PQ}) = w (\mathrm{P}) + w (\mathrm{Q})\tag{3}
\]

对于 \(\mathbf{P} = \sum_{j} a_{j} \mathbf{X}^{j}\) 和 \(\mathbf{Q} = \sum_{j} b_{j} \mathbf{X}^{j}\)。令 \(i\)（分别地，\(k\)）为下述整数 \(j\) 中的最小者：它们使 \(v(a_{j}) + j\xi\)（分别地，\(v(b_{j}) + j\xi\)）取得最小值；令 \(\alpha\)（分别地，\(\beta\)）表示这个最小值。对于 \(j, j'\) 属于 \(\mathbf{N}\)，

\[
w (a _ {j} b _ {j ^ {\prime}} \mathbf {X} ^ {j + j ^ {\prime}}) = v (a _ {j}) + j \xi + v (b _ {j ^ {\prime}}) + j ^ {\prime} \xi \geqslant \alpha + \beta ,
\]

由此根据 (2) 得 \(w(\mathrm{PQ}) \geqslant \alpha + \beta\)。现在考虑项 \(c\mathbf{X}^{i + k}\)，它在 PQ 中的次数为 \(i + k\)；则 \(c = \sum_{n \in \mathbf{Z}} a_{i + n} b_{k - n}\)；根据 \(i\) 和 \(k\) 的选取，该元素

\[
w (a _ {i + n} b _ {k - n} \mathbf {X} ^ {i + k}) = v (a _ {i + n}) + (i + n) \xi + v (b _ {k - n}) + (k - n) \xi
\]

其最小值 \(\alpha + \beta\) 仅在 \(n = 0\) 时且只在此时取得；因此 \(w(c\mathbf{X}^{i + k}) = \alpha + \beta\)，从而根据 (1)，

\[
w (\mathrm{PQ}) = \alpha + \beta = w (\mathrm{P}) + w (\mathrm{Q}).
\]

命题 1. 设 K 为域，v 为 K 上的赋值，\(\Gamma\) 为其阶群，\(\Gamma'\) 为包含 \(\Gamma\) 的全序群，且 \(\xi\) 是 \(\Gamma'\) 的一个元素，并满足关系 \(n\xi\in\Gamma\)，\(n\in Z\) 蕴含 n=0。则存在唯一的赋值 w，取值于 \(\Gamma'\) 且延拓 v，并满足 \(w(\mathbf{X})=\xi\)。w 的剩余域等于 v 的剩余域，其阶群是 \(\Gamma+Z\xi\)（\(\Gamma'\) 的子群）。

我们先证明 \(w\) 的唯一性。设 \(\mathbf{P} = \sum_{j} a_j \mathbf{X}^j\) 是 \(\mathbf{K}[\mathbf{X}]\) 中的元素。则 \(w(a_j \mathbf{X}^j) = v(a_j) + j\xi\)，这说明单项式 \(a_j \mathbf{X}^j\) 在 \(a_j \neq 0\) 时关于 \(w\) 的值彼此不同。由此 \(w(\mathbf{P}) = \inf_j (v(a_j) + j\xi)\)，这同时说明 \(w\) 在 \(\mathbf{K}[\mathbf{X}]\) 上（从而也在 \(\mathbf{K}(\mathbf{X})\) 上）的唯一性，并说明 \(w\) 的阶群是 \(\Gamma + \mathbf{Z}\xi\)。还可看出，若 \(\mathbf{P} \neq 0\)，则可写成 \(\mathbf{P} = a\mathbf{X}^n (1 + u)\) 的形式，其中 \(a\in \mathbf{K}^*\)，\(n\in \mathbf{N}\)，\(u\in \mathbf{K}(\mathbf{X})\)，且 \(w(u) > 0\)；因此，对于任意 \(\mathbf{R}\neq 0\)，它作为 \(\mathbf{K}(\mathbf{X})\) 中的元素都可写成

\[
\mathrm{R} = b \mathrm{X} ^ {n} (1 + u ^ {\prime}),
\]

其中 \(b \in \mathbf{K}^*\)，\(n \in \mathbf{Z}\)，\(u' \in \mathbf{K}(\mathbf{X})\)，且 \(w(u') > 0\)；于是 \(w(\mathbf{R}) = v(b) + n\xi\)，所以 \(w(\mathbf{R}) = 0\) 当且仅当 \(v(b) = 0\) 且 \(n = 0\)；因此，当 \(w(\mathbf{R}) = 0\) 时，\(\mathbf{R}\) 与 \(b\) 模 \(w\) 的理想同余，这说明 \(w\) 的剩余域等于 \(v\) 的剩余域。

最后，w 的存在性由引理 1 得出。

命题 2. 设 K 为域，v 为 K 上的赋值，\(\Gamma\) 为其阶群，k 为其剩余域。存在 K(X) 上唯一的延拓 v 的赋值 w，使得 \(w(X) = 0\)，且 X 在剩余域 \(k'\)（w 的剩余域）中的像 t 超越于 k。w 的阶群等于 v 的阶群，其剩余域为 \(k(t)\)。

为证明 \(w\) 的唯一性，只需证明：若 \(\mathbf{P} = \sum_{j} a_j \mathbf{X}^j\) 是 \(\mathbf{K}[\mathbf{X}]\) 中的非零元素，则

\[
w (\mathbf {P}) = \inf _ {j} (v (a _ {j})).
\]

可以用 \(K^{*}\) 中的一个元素去除 P，并设对所有 j 都有 \(v(a_{j}) \geqslant 0\)，且其中一个 \(v(a_{j})\) 为 0。由于 \(w(\mathbf{X}) = 0\)，P 于是属于 w 的环；记 \(\bar{a}_{j}\) 为 \(a_{j}\) 在 k 中的典范像，则 P 在剩余域 \(k'\) 中的典范像为 \(\sum_{j} \bar{a}_{j} t^{j}\)；由于 t 超越于 k 且其中一个 \(\bar{a}_{j}\) 非零，该像非零，从而

\[
w (\mathbf {P}) = 0 = \inf _ {j} (v (a _ {j})).
\]

现在证明 \(w\) 的存在性。根据引理 1，公式 \(w(\mathbf{P}) = \inf_{j}(v(a_j))\)（对于 \(\mathbf{P} = \sum_{j} a_j \mathbf{X}^j\)）定义赋值 \(w\)，作用在 \(\mathbf{K}(\mathbf{X})\) 上，且显然 \(w\) 与 \(v\) 有相同的阶群。于是 \(w(\mathbf{X}) = 0\)。我们证明 \(t\)（\(\mathbf{X}\) 在剩余域 \(k'\)（\(w\) 的剩余域）中的典范像）超越于 \(k\)：若 \(\sum_{j} \bar{a}_j t^j = 0\)，其中 \(\bar{a}_j \in k\) 对所有 \(j\) 成立，则取 \(a_j\) 为 \(\bar{a}_j\) 在 \(v\) 的环中的代表，有 \(w\left(\sum_{j} a_j \mathbf{X}^j\right) > 0\)；从而 \(v(a_j) > 0\) 对所有 \(j\) 成立，于是 \(\bar{a}_j = 0\) 对所有 \(j\) 成立。最后证明 \(k' = k(t)\)：每个元素 \(\mathbf{R}\) 都属于 \(\mathbf{K}(\mathbf{X})\)，且都可写成 \(\mathbf{R} = c\left(\sum_{j} a_j \mathbf{X}^j\right) / \left(\sum_{j} b_j \mathbf{X}^j\right)\) 的形式，其中 \(c, a_j, b_j\) 属于 \(\mathbf{K}\)，\(v(a_j) \geqslant 0\) 且 \(v(b_j) \geqslant 0\) 对所有 \(j\) 成立，其中一个 \(v(a_j)\) 和一个 \(v(b_j)\) 为零；于是 \(w(\mathbf{R}) \geqslant 0\) 当且仅当 \(v(c) \geqslant 0\)；记 \(f\) 为从 \(w\) 的环到 \(k'\) 的典范同态，

\[
f (\mathrm{R}) = f (c) \left(\sum_ {j} f (a _ {j}) t ^ {j}\right) / \left(\sum_ {j} f (b _ {j}) t ^ {j}\right),
\]

这就证明了我们的断言。

注。不要以为我们刚遇到的 v 到 K(X) 的两类延拓是仅有的两类；还可能存在第三类延拓，其中 \(\Gamma'/\Gamma\) 是扭群，而 \(k'\) 是 k 的（不必有限的）代数扩张。这第三类不一定由引理 1 中所述的步骤得到（参见~§3，习题 1）。

